\vspace*{-1em}\section{Introduction}
\label{sec:intro}

The camera is the visual portal to the filmmaker's world, guiding the audience's gaze as the story unfolds and constructing the narrative's visual language. While traditional visual effects can dramatically transform a raw film set into an immersive spectacle, the ability to manipulate the camera during post-production introduces another dimension of control over visual storytelling.

To this end, we synthesize or `render' the dynamic scene specified by an input source video from novel camera trajectories and viewpoints, which we call \emph{video reshooting}. Importantly, we must achieve faithful reconstruction of seen content in the source video and photorealistically plausible generation of unseen content, all while maintaining precise, user-definable camera control.

We will employ video diffusion models since they are powerful priors for generating dynamic content which is geometrically and temporally coherent \cite{brooks2024sora, wan2025, wiedemer2025videozeroshot, kong2025hunyuanvideosystematicframeworklarge, nvidia2025worldsimulationvideofoundation, genmo2024mochi}.  We will further combine the diffusion models with 4D reconstruction which lifts the monocular source video into a 4D point cloud, providing spatiotemporal grounding for reconstruction and a rich signal for camera control. We present \emph{\methodname}, a video reshooting framework that grounds the source video and target cameras in a 4D point cloud with temporally-persistent static pixels, while leveraging the generative priors of video diffusion models.

Existing works for video reshooting \cite{yu2025trajcraft, ren2025gen3c, hu2025ex4d} condition video diffusion models on per-frame depth-lifted point clouds rendered in the target cameras. However, they often struggle with geometry artifacts and/or temporal flickering due to imprecise 4D reconstruction of real-world dynamic videos as they are often trained on point cloud renders from precise depth maps. Moreover, they also struggle with accurate camera control and content preservation with challenging target camera trajectories and viewpoints.

\methodname introduces the following key designs that not only show state-of-the-art visual quality and robustness to a wide variety of source videos and target cameras but also extend our method with capabilities beyond vanilla video reshooting. First, we build a 4D-grounded point cloud representation where static pixels are visible from any frame via segmentation and 4D reconstruction, as opposed to the per-frame 3D point cloud of baselines. Conditioning with static pixel temporal persistence establishes both explicit preservation of seen content and provides rich camera signals even when the target cameras have little per-frame overlap with the source video. Second, we augment model training with dynamic, 4D-reconstructed multiview video pairs that contain depth estimation artifacts from non-frontal views. Thus, \methodname is significantly more robust to the quality of real-world point cloud renders while allowing us to additionally condition on the source video to utilize video model priors for geometric coherence. This further enables us to manipulate the 4D point cloud during inference for real-world applications beyond video reshooting.

Our contributions are as follows:
\begin{itemize}
    \item We present \methodname, a video reshooting framework that maintains geometric and physical plausibility with real-world inference, while explicitly preserving seen content by grounding generation in a 4D point cloud.
    \item Through extensive quantitative and qualitative comparisons, including a user study, we validate the improved content preservation, camera controllability, and visual quality of \methodname over state-of-the-art baselines for a wide variety of videos and cameras.
    \item We show that our training extends \methodname with capabilities that generalize to real-world applications such as dynamic scene expansion, 4D scene recomposition, and long video inference with memory.
\end{itemize}
